\documentclass[UTF8,11pt]{ctexart}
\usepackage[a4paper,margin=24mm]{geometry}
\usepackage{amsmath,amssymb,amsthm,mathtools,mathrsfs}
\usepackage[all]{xy}
\usepackage{xurl,hyperref,enumitem}
\usepackage{tikz}
\usetikzlibrary{arrows.meta,cd,decorations.markings,calc,patterns}
\hypersetup{hidelinks,breaklinks=true}
\setlength{\emergencystretch}{3em}
\allowdisplaybreaks
\newtheorem*{theorem}{Theorem}
\newtheorem*{thm}{Theorem}
\newtheorem*{lemma}{Lemma}
\newtheorem*{proposition}{Proposition}
\newtheorem*{prop}{Proposition}
\newtheorem*{corollary}{Corollary}
\newtheorem*{cor}{Corollary}
\newtheorem*{claim}{Claim}
\theoremstyle{definition}
\newtheorem*{definition}{Definition}
\newtheorem*{defn}{Definition}
\newtheorem*{remark}{Remark}
\newtheorem*{example}{Example}
\newcommand{\R}{\mathbb R}
\newcommand{\C}{\mathbb C}
\newcommand{\Z}{\mathbb Z}
\newcommand{\Q}{\mathbb Q}
\newcommand{\N}{\mathbb N}
\newcommand{\F}{\mathbb F}
\newcommand{\D}{\mathbb D}
\newcommand{\bB}{\mathbb B}
\newcommand{\bC}{\mathbb C}
\newcommand{\bR}{\mathbb R}
\newcommand{\bZ}{\mathbb Z}
\newcommand{\bN}{\mathbb N}
\newcommand{\bQ}{\mathbb Q}
\newcommand{\bD}{\mathbb D}
\newcommand{\bF}{\mathbb F}
\newcommand{\CP}{\mathbb{CP}}
\newcommand{\RP}{\mathbb{RP}}
\newcommand{\sO}{\mathscr O}
\newcommand{\sI}{\mathscr I}
\newcommand{\sF}{\mathscr F}
\newcommand{\sG}{\mathscr G}
\newcommand{\sH}{\mathscr H}
\newcommand{\sR}{\mathscr R}
\newcommand{\sM}{\mathscr M}
\newcommand{\sV}{\mathscr V}
\newcommand{\sU}{\mathscr U}
\DeclareMathOperator{\Hom}{Hom}
\DeclareMathOperator{\Ext}{Ext}
\DeclareMathOperator{\Tor}{Tor}
\DeclareMathOperator{\Spec}{Spec}
\DeclareMathOperator{\Proj}{Proj}
\DeclareMathOperator{\supp}{supp}
\DeclareMathOperator{\im}{im}
\DeclareMathOperator{\id}{id}
\DeclareMathOperator{\Tr}{Tr}
\DeclareMathOperator{\tr}{tr}
\DeclareMathOperator{\rank}{rank}
\DeclareMathOperator{\ord}{ord}
\DeclareMathOperator{\dist}{dist}
\DeclareMathOperator{\End}{End}
\DeclareMathOperator{\Aut}{Aut}
\DeclareMathOperator{\Gal}{Gal}
\DeclareMathOperator{\vol}{vol}
\DeclareMathOperator{\Cl}{Cl}
\DeclareMathOperator{\coker}{coker}
\DeclareMathOperator{\codim}{codim}
\DeclareMathOperator{\divergence}{div}
\DeclareMathOperator{\Ric}{Ric}
\DeclareMathOperator{\grad}{grad}
\DeclareMathOperator{\Hess}{Hess}
\newcommand{\abs}[1]{\left\lvert#1\right\rvert}
\newcommand{\sabs}[1]{\lvert#1\rvert}
\newcommand{\babs}[1]{\bigl\lvert#1\bigr\rvert}
\newcommand{\Babs}[1]{\Bigl\lvert#1\Bigr\rvert}
\newcommand{\bbabs}[1]{\biggl\lvert#1\biggr\rvert}
\newcommand{\BBabs}[1]{\Biggl\lvert#1\Biggr\rvert}
\newcommand{\norm}[1]{\left\lVert#1\right\rVert}
\newcommand{\snorm}[1]{\lVert#1\rVert}
\newcommand{\bnorm}[1]{\bigl\lVert#1\bigr\rVert}
\newcommand{\Bnorm}[1]{\Bigl\lVert#1\Bigr\rVert}
\newcommand{\bbnorm}[1]{\biggl\lVert#1\biggr\rVert}
\newcommand{\BBnorm}[1]{\Biggl\lVert#1\Biggr\rVert}
\newcommand{\widebar}[1]{\overline{#1}}
\newcommand{\nicefrac}[2]{\frac{#1}{#2}}
\newcommand{\linnprod}[2]{\langle#1,#2\rangle}
\newcommand{\myindex}[1]{#1}
\newcommand{\glsadd}[1]{}
\newcommand{\avoidbreak}{}
\newcommand{\sourceref}[1]{\textnormal{[原文：\nolinkurl{#1}]}}
\newcommand{\thmref}[1]{\sourceref{#1}}
\newcommand{\lemmaref}[1]{\sourceref{#1}}
\newcommand{\propref}[1]{\sourceref{#1}}
\newcommand{\corref}[1]{\sourceref{#1}}
\newcommand{\defnref}[1]{\sourceref{#1}}
\newcommand{\exampleref}[1]{\sourceref{#1}}
\newcommand{\exerciseref}[1]{\sourceref{#1}}
\newcommand{\figureref}[1]{\sourceref{#1}}
\newcommand{\sectionref}[1]{\sourceref{#1}}
\newcommand{\appendixref}[1]{\sourceref{#1}}
\newcommand{\stacksref}[2]{\href{https://stacks.math.columbia.edu/tag/#1}{#2 [#1]}}


\newcommand{\E}{\mathbb E}
\newcommand{\Prob}{\mathbb P}
\newcommand{\Reff}{\mathcal R_{\mathrm{eff}}}
\newcommand{\eps}{\varepsilon}
\DeclareMathOperator{\diam}{diam}
\DeclareMathOperator{\Var}{Var}
\DeclareMathOperator{\Crit}{Crit}
\newcommand{\dd}{\,\mathrm d}

\begin{document}
\section*{051\quad 均匀生成森林的自由与连线极限}
\subsection*{来源与计数目标}
Asaf Nachmias, \emph{Planar Maps, Random Walks and Circle Packing}, Lecture Notes in Mathematics 2243, Springer, 2020（在线出版 2019）。第7章 \emph{Uniform Spanning Trees of Planar Graphs}，\S7.2--7.3，Theorems 7.2--7.3、Claim 7.4 及空间 Markov 性，章页89--103中相应小节。实际读取 Springer 开放获取 HTML 数学正文，获取日期2026-09-28。
\par\url{https://link.springer.com/chapter/10.1007/978-3-030-27968-4_7}
\par\textbf{本文件只计一个问题：}证明有限耗尽上的均匀生成树在自由和连线边界条件下均有与耗尽无关的柱集概率极限（Theorem 7.3）。
\subsection*{作者命题与人类证明}
\paragraph{Setting (Sections 7.1 and 7.3).}
\noindent\textit{以下背景与记号据所列原文章节摘编；不是逐字引文，也不是新增证明。}\par
Let $G$ be a finite connected graph. A spanning tree $T$ of $G$ is a connected subgraph of $G$ that contains no cycles and such that every vertex of $G$ is incident to at least one edge of $T$. The set of spanning trees of a given finite connected graph is obviously finite and hence we may draw one uniformly at random. This random tree is called the uniform spanning tree (UST) of $G$.

Let $G$ be an infinite connected graph and let $\{G_n\}$ be a finite exhaustion of it. Denote by $G_n^*$ the graph obtained from $G$ by identifying the infinite set of vertices $G\setminus G_n$ to a single vertex $z_n$ and erasing the loops at $z_n$ formed by this identification. We say that $\{G_n^*\}$ is a wired finite exhaustion of $G$.

\begin{theorem}[7.3, Pemantle]
Let $G$ be an infinite connected graph, $\{G_n\}$ a finite exhaustion and $\{G_n^*\}$ the corresponding wired finite exhaustion. Denote by $\mathcal T_n$ and $\mathcal T_n^*$ USTs on $G_n$ and $G_n^*$, respectively. Then for any two finite disjoint subsets $A,B\subset E(G)$ of edges of $G$ we have that the limits
\[
\lim_{n\to\infty}\Prob(A\subset\mathcal T_n,\ B\cap\mathcal T_n=\emptyset)
\quad\text{and}\quad
\lim_{n\to\infty}\Prob(A\subset\mathcal T_n^*,\ B\cap\mathcal T_n^*=\emptyset)
\]
exist and do not depend on the exhaustion $\{G_n\}$.
\end{theorem}

\paragraph{Supporting argument: Kirchhoff's effective resistance formula (Section 7.2.1).}
\begin{theorem}[7.2, Kirchoff]
Let $G$ be a finite connected graph and denote by $\mathcal T$ a uniformly drawn spanning tree of $G$. Then for any edge $e=(x,y)$ we have
\[
\Prob(e\in\mathcal T)=\Reff(x\leftrightarrow y).
\]
\end{theorem}
\begin{proof}
Let $a\ne z$ be two distinct vertices of $G$ (later we will take $a=x$ and $z=y$) and note that any spanning tree of $G$ contains precisely one path connecting $a$ and $z$. Thus, a uniformly drawn spanning tree induces a random path from $a$ to $z$. By Claim 2.46 we obtain a unit flow $\theta$ from $a$ to $z$.
To be concrete, for each edge $e$ we have that $\theta(\vec e)$ is the probability that the random path from $a$ to $z$ traverses $\vec e$ minus the probability that it traverses $\overleftarrow e$. We will now show that $\theta$ satisfies the cycle law (see Claim 2.14), so it is in fact the unit current flow (see Definition 2.19).

Let $\vec e_1,\ldots,\vec e_m$ be a directed cycle in $G$. Our goal is to show that
\[
\sum_{i=1}^m\theta(\vec e_i)=0.\tag{7.2}
\]
Denote by $T(G)$ the set of spanning trees of $G$. Expanding the sum on the left hand side with the definition of $\theta$ we get that it equals
\[
|T(G)|^{-1}\sum_{t\in T(G)}\sum_{i=1}^m f_i^+(t)
-|T(G)|^{-1}\sum_{t\in T(G)}\sum_{j=1}^m f_j^-(t),
\]
where $f_i^+(t)$ equals $1$ if the unique path from $a$ to $z$ in $t$ traverses $\vec e_i$ and $0$ otherwise, and similarly, $f_j^-(t)$ equals $1$ if this path traverses $\overleftarrow e_j$ and $0$ otherwise.

For $1\le i\le m$ we denote by $T_i^+$ the set of pairs $(t,i)$ for which $f_i^+(t)=1$. Similarly define $T_j^-$ as the set of pairs $(t,j)$ for which $f_j^-(t)=1$. To prove (7.2) it suffices to show that
\[
\left|\biguplus_{i\in\{1,\ldots,m\}}T_i^+\right|
=\left|\biguplus_{j\in\{1,\ldots,m\}}T_j^-\right|.
\]
Let $(t,i)\in T_i^+$. The graph $t\setminus\{e_i\}$ has two connected components. Let $\vec e_j$ be the first edge after $\vec e_i$, in the order of the cycle $\vec e_1,\ldots,\vec e_m$, that is incident to both connected components and consider the spanning tree $t'=t\cup\{e_j\}\setminus\{e_i\}$. Note that the unique path in $t'$ from $a$ to $z$ traverses $\overleftarrow e_j$, so $(t',j)\in T_j^-$. This procedure defines a bijection from $\biguplus_i T_i^+$ to $\biguplus_j T_j^-$.
Indeed, given $(t',j)$ from before, we can erase $e_j$ and go on the cycle in the opposite order until we reach $e_i$ which has to be the first edge incident to the two connected components of $t'\setminus\{e_j\}$. This shows (7.2) and concludes the proof.
\end{proof}

\paragraph{Spatial Markov property (Section 7.2.2).}
Let $G=(V,E)$ be a finite connected graph and let $A$ and $B$ be two disjoint subsets of edges. We write $(G-B)/A$ for the graph obtained from $G$ by erasing the edges of $B$ and contracting the edges of $A$. We identify the edges of $(G-B)/A$ with the edges $E\setminus B$. Denote by $\mathcal T_G$ and $\mathcal T_{(G-B)/A}$ a UST on $G$ and $(G-B)/A$, respectively, and assume that
\[
\Prob(A\subset\mathcal T_G,\ B\cap\mathcal T_G=\emptyset)>0.
\]
This assumption is equivalent to $G-B$ being connected and that $A$ contains no cycles.
Then, conditioned on the event that $\mathcal T_G$ contains the edges $A$ and does not contain any edge of $B$ the distribution of $\mathcal T_G$ is equal to the union of $A$ with $\mathcal T_{(G-B)/A}$. In other words, for a set $\mathcal A$ of spanning trees of $G$ we have that
\begin{multline*}
\Prob(\mathcal T_G\in\mathcal A\mid A\subset\mathcal T_G,\ B\cap\mathcal T_G=\emptyset)\\
=\Prob(A\cup\mathcal T_{(G-B)/A}\in\mathcal A).\tag{7.3}
\end{multline*}
\begin{proof}
The proof of (7.3) follows immediately from the observation that the set of spanning trees of $G$ not containing any edge of $B$ is simply the set of spanning trees of $G-B$. Similarly, the set of spanning trees of $G$ containing all the edges of $A$ is simply the union of $A$ to each spanning tree of $G/A$, and (7.3) follows.
\end{proof}

\begin{claim}[7.4]
Let $G$ be an infinite connected graph, $\{G_n\}$ a finite exhaustion and $\{G_n^*\}$ a wired finite exhaustion. Then for any two vertices $u,v$ of $G$ we have that the limits
\[
\Reff^F(u\leftrightarrow v;G):=\lim_n\Reff(u\leftrightarrow v;G_n),
\qquad
\Reff^W(u\leftrightarrow v;G):=\lim_n\Reff(u\leftrightarrow v;G_n^*)
\]
exist and do not depend on the exhaustion $\{G_n\}$.
\end{claim}
\begin{proof}
For the first limit we note that by Rayleigh's monotonicity (Corollary 2.29), the sequence $\Reff(u\leftrightarrow v;G_n)$ is non-increasing and non-negative since $G_n\subset G_{n+1}$, hence it converges. A sandwiching argument as in the proof of Claim 7.4 shows that the limit does not depend on the exhaustion $\{G_n\}$.

For the second limit, since $G_n$ can be obtained by gluing vertices of $G_{n+1}$ we deduce by Corollary 2.30 that the sequence $\Reff(u\leftrightarrow v;G_n^*)$ is non-decreasing and bounded (by the graph distance in $G$ between $u$ and $v$ for instance), hence it converges. The limit does not depend on the exhaustion by an identical sandwiching argument.
\end{proof}

\begin{proof}[Proof of Theorem 7.3 (Section 7.3.1)]
We will prove the assertion regarding the first limit; the second is almost identical. Write $A=\{e_1,\ldots,e_k\}$ and $e_i=(x_i,y_i)$ for each $1\le i\le k$. Assume without loss of generality that $G_n$ contains $A$ for all $n$. As before, denote by $\mathcal T_n$ a UST of $G_n$. By (7.3) and Theorem 7.2 we have that
\begin{align*}
\Prob(A\subset\mathcal T_n)
&=\prod_{i=1}^k\Prob(e_i\in\mathcal T_n\mid e_j\in\mathcal T_n\quad\forall j<i)\\
&=\prod_{i=1}^k\Reff(x_i\leftrightarrow y_i;G_n/\{e_1,\ldots,e_{i-1}\}).
\end{align*}
Note that $\{G_n/\{e_1,\ldots,e_{i-1}\}\}$ is a finite exhaustion of the infinite graph $G/\{e_1,\ldots,e_{i-1}\}$ and so by Claim 7.4 we obtain that the limit
\[
\lim_n\Prob(A\subset\mathcal T_n)
=\prod_{i=1}^k\Reff(x_i\leftrightarrow y_i;G/\{e_1,\ldots,e_{i-1}\}),
\]
exists and does not depend on the exhaustion.
Since we know this limit exists for all finite edge sets $A$, it follows by the inclusion-exclusion formula that $\Prob(A\subset\mathcal T_n,B\cap\mathcal T_n=\emptyset)$ converges for any finite sets $A,B$, concluding our proof.
\end{proof}
\subsection*{整理说明（非作者正文）}
先把主问题列明，再按作者次序收入其核心支撑和最终证明；不另计 Kirchhoff 公式或电阻极限。使用范围沿全书为局部有限图，边电导为1；$\overleftarrow e$ 表示与 $\vec e$ 相反的有向边。先修为电流的节点律和循环律、Rayleigh 单调性、条件概率乘法公式及容斥原理。难点在于生成树换边的计数双射、条件分布与图收缩的对应，以及把这些联系传递到耗尽极限。

原文 Theorem 7.3 明说只详细写自由情形，连线情形用同一论证；两个边界电阻极限的论证均已收录。Claim 7.4 原文有同号自引及第二段 $G_n,G_{n+1}$ 未印星号，照录并在此指出，不将编辑修补混入原文。空间 Markov 段的边集识别沿原句保留。主证明未逐列处理零概率条件事件；本条保留作者的陈述范围和证明写法，不另造补证。符号、换行和超链接引用转为独立 TeX；这是网页转录，不是下载的原生 TeX 或 PDF。
\subsection*{可修改的简短标注（非作者正文）}
$c$：无限图上随机生成森林概率测度的构造。

$a$：生成树换边双射；路径诱导单位流；有效电阻；图收缩；空间 Markov 性；Rayleigh 单调性；容斥原理。

\texttt{cross\_theory\_status = yes}。生成树中的边事件经换边双射和电网络公式转为有效电阻，条件化转为图收缩，再由电阻单调极限得到概率极限。位置：Theorem 7.2 与 \S7.3.1。
标签为非穷尽、可修改初稿。
\subsection*{许可与修改记录}
原作者及作品见来源栏；来源采用 CC BY 4.0。本题证摘录和附加整理沿用该许可。2026-09-28：增加题号、来源定位、独立排版、整理说明和初稿标注；数学内容的取材范围及排版变化见整理说明。
\end{document}
